\scanpage{113}
\mspage{58}
% archive sheet number: 112

\textcircled{\scriptsize\unsure{D$_4$|2}} Les localisés pour les $W$-\uline{équivalences fines} \uline{droites}.

Il n'y a plus lieu, on l'a dit, de regarder les catégories \unsure{internes|externes|extrêmes} du diagramme p.~8, où ladite équivalence \deleted{\unsure{coï}}coïncide avec la $W$-équiv. grossière, pour laquelle on connaît déjà les localisés. Il reste à regarder les catégories

\[
 \deleted{\unsure{\mathrm{Fib}}}\mathrm{Cat}_{/S},\quad \underline{\mathrm{Liss}}\,S,\quad \underline{\mathrm{Fib}}\,S,\quad \underline{\mathrm{Fibsc}}\,S .
\]

Le lemme de la page 11 nous dit \deleted{\unsure{que}} qu'il suffit de montrer que dans la situation

\[
 \underline{\mathrm{Fibsc}}\,S\xrightarrow{\lambda}\mathrm{Cat}_{/S},
\]

il y a un foncteur en sens inverse, avec les propriétés spécifiées. Il en résultera que toute sous-catégorie pleine $\unsure{F|\mathcal F}$ de $\mathrm{Cat}_{/S}$ qui contient $\underline{\mathrm{Fib}}\,S$, quand on la munit du localiseur de la $W$-équivalence fine droite, donne un localisé équivalent, par \deleted{\unsure{le}}\unsure{le} foncteur induit par l'inclusion, au localisé correspondant de $\mathrm{Cat}_{/S}$, i.e. essentiellement : $\mathrm{HOT}_W(S)$.

Comme foncteur en sens inverse, il s'impose de prendre encore $\Phi$. Il y a à vérifier trois choses :

a) $\Phi$ compatible avec les localiseurs, i.e si $f\colon$

\diagram{\begin{tikzcd}[ampersand replacement=\&,row sep=tiny,column sep=tiny]
 X\arrow[rr]\arrow[dr,dash]\&\&Y\arrow[dl,dash]\\
 \&S\&
\end{tikzcd}}

sur $S$ est une $W$-équiv. fine droite, i.e. les $ {}_{s\backslash}X\mathbin{\unsure{-|\to}}{}_{s\backslash}Y$ sont des $W$-équivalences, alors $\Phi(f)$ aussi. Mais la \unsure{dite} équivalence coïncide

